\section{Marked lines and the rank inequality}\label{sec:lines}

The finite-field construction supplies two features at once: a sparse
incidence pattern with large girth, and many matrices whose total rank
is bounded by a much smaller sum of local ranks. The metric in the next
section will use the girth. The rank argument uses only the marked
points and the triangular relations.
Lemma~\ref{lem:incidence} adapts \cite[Lemma~2.2]{OAI252}, which is
stated in dimension \(20\) with girth greater than \(12\), to dimension
\(d=103\) and girth greater than \(g=100\).

\subsection{Incidence data with explicit parameters}

Fix
\begin{equation}\label{eq:constants}
 g=100,\qquad d=103,\qquad
 c=(20d)^{-d}=2060^{-103},\qquad r=100c^{-1}.
\end{equation}
The integer \(g\) is the required girth bound and \(d\) is the affine
dimension. Let \(p\) be a prime, and let \(Q=p^a\) be any power of
\(p\) greater than
\begin{equation}\label{eq:qbound}
 Q_0=320000\cdot2060^{309}.
\end{equation}
This is the largest of
\[
 20d,\qquad 32r^2/c,\qquad 256r/c,\qquad 16rg^2.
\]
Fix a monic irreducible polynomial \(\pi\in\F_p[X]\) of degree \(a\),
and put \(\F_Q=\F_p[X]/(\pi)\). Each element of \(\F_Q\) is uniquely
\(c_0+c_1X+\cdots+c_{a-1}X^{a-1}\) with residues
\(c_j\in\{0,\ldots,p-1\}\); give it the number
\(c_0+c_1p+\cdots+c_{a-1}p^{a-1}\in\{0,\ldots,Q-1\}\).
This numbering orders \(\F_Q\); it plays no role in its arithmetic.
Put \(P=\F_Q^d\), \(m=Q^d\), and let \(S\subset\F_Q\) consist of the
elements numbered
\[
 0,1,\ldots,\lfloor Q/(10d)\rfloor-1.
\]
Only \(|S|=\lfloor Q/(10d)\rfloor\) matters below.
An affine line \(\ell_i\) has a distinguished mark \(a_i\in S^d\).
Its \emph{punctured line} is \(\ell_i^\circ=\ell_i\setminus\{a_i\}\).
Marks, and hence labels \(i\), will be distinct. Different labels
are allowed initially to choose the same geometric line.

\begin{lemma}\label{lem:incidence}
There are a nonempty finite label set \(I\), marked affine lines
\((\ell_i,a_i)_{i\in I}\), and a set \(\cE\subset P\) such that,
with \(N=|I|\),
\begin{enumerate}
\item the bipartite incidence graph on \(P\sqcup I\), with
\(z\sim i\) when \(z\in\ell_i^\circ\), has no cycle of length at most \(g\);
\item every \(z\notin\cE\) belongs to at least \(4r^2\) punctured lines;
\item \(N\ge cm/2\) and \(|\cE|\le m/(2r)\).
\end{enumerate}
Consequently,
\begin{equation}\label{eq:budget}
 2m+r|\cE|\le\frac{5m}{2}
 <\frac{25m}{2}\le\frac{Nr}{4}.
\end{equation}
\end{lemma}

\begin{proof}
Initially use all \(n_0=|S|^d\) points of \(S^d\) as marks. At each
mark choose a direction independently and uniformly from the
\[
 J=\frac{Q^d-1}{Q-1}
\]
projective directions. Since \(Q>20d\), we have
\(|S|\ge Q/(20d)\), hence \(n_0\ge cm\).
Also \(Q^{d-1}\le J\le2Q^{d-1}\).

For a fixed \(z\), each mark other than \(z\) produces an incidence
with probability \(1/J\); the mark at \(z\), if present, produces none.
Thus the degree is binomial, with mean \(\mu_z\) and variance at most
\(\mu_z\). The bounds on \(Q\) imply \(cm>2\), so
\[
 \mu_z\ge\frac{n_0-1}{2Q^{d-1}}
 \ge\frac{cQ}{4}>8r^2.
\]
Chebyshev's inequality gives
\[
 \Pr\{\deg(z)<4r^2\}\le\frac{4}{\mu_z}
 \le\frac{16}{cQ}.
\]
The expected number of points of degree less than \(4r^2\) is therefore
at most \(16m/(cQ)\). By Markov's inequality, the probability that
there are more than \(m/(4r)\) such points is at most
\[
 \frac{64r}{cQ}<\frac14.
\]

Next count simple cycles of length \(2s\), where \(2\le s\le g/2\).
There are at most \(m^s\) ordered choices of their point vertices.
Consecutive points determine a geometric line, which has at most \(Q\)
possible marks. If the \(s\) chosen marks are distinct, the required
directions occur with probability \(J^{-s}\).
Ignoring distinctness restrictions and forbidden punctured endpoints
only increases the count. Thus the expected number of such cycles is
at most
\[
 m^s Q^sJ^{-s}\le Q^{2s}.
\]
The expected total number of cycles of length at most \(g\) is at most
\(gQ^g\). With probability at least \(3/4\), their number is at most
\(4gQ^g\). Both desired bounds therefore hold simultaneously with
positive probability.

Choose such an outcome. Delete every label on a cycle of length at
most \(g\). At most \(4g^2Q^g\) labels are deleted. Let \(\cE\)
contain the old low-degree points and all points on deleted lines.
Deletion creates no new cycle, and changes no degree outside \(\cE\).
Because \(d=g+3\),
\[
 \frac{\#\{\text{deleted labels}\}}m\le\frac{4g^2}{Q^3},
 \qquad
 \frac{\#\{\text{points on deleted lines}\}}m
 \le\frac{4g^2}{Q^2}.
\]
The inequality \(Q>16rg^2\), together with \(r=100/c\), makes the first
quantity at most \(c/2\) and the second at most \(1/(4r)\).
It follows that \(N\ge cm/2\) and
\(|\cE|\le m/(4r)+m/(4r)\). Finally
\[
 Nr/4\ge cmr/8=25m/2,
\]
which proves~\eqref{eq:budget}.
\end{proof}

The probabilistic proof makes the data finite and deterministic as
follows. Normalize each nonzero direction vector by making its first
nonzero coordinate \(1\). Order these vectors and the marks
lexicographically, comparing coordinates by their numbers. Choose the first direction
tuple meeting the two bounds before deletion, then delete exactly the
labels described in the proof. Lemma~\ref{lem:incidence} ensures that
this finite search terminates.

\subsection{Two blocks at each regular point}

Write
\[
 I(z)=\{i\in I:z\in\ell_i^\circ\}.
\]
For each regular point \(z\notin\cE\), partition \(I(z)\) into two
blocks \(I_b\), each with
\begin{equation}\label{eq:block-shape}
 |I_b|=s_bt_b,\qquad s_b,t_b\ge r.
\end{equation}
Here is a definite choice. If \(n=|I(z)|\ge4r^2\), put
\[
 b_n=r+(n\bmod r),\qquad
 a_n=\frac{n-(r+1)b_n}{r}.
\]
Then \(r\le b_n\le2r-1\), the number \(a_n\) is integral, and
\[
 a_n\ge\frac{2r^2-r+1}{r}\ge r.
\]
Use a first block with \((s_b,t_b)=(r,a_n)\) and a second with
\((s_b,t_b)=(r+1,b_n)\).
Order the labels and use the initial segment for the first block.
There are no blocks at points of \(\cE\).

Let \(\cB\) be the set of blocks, and let \(z_b\) be the point at
which block \(b\) is based. Fix the lexicographic bijection
\begin{equation}\label{eq:coordinates}
 I_b\longrightarrow\Z/s_b\Z\times\Z/t_b\Z,\qquad
 i\longmapsto(\alpha_b(i),\beta_b(i)).
\end{equation}
The block--line graph on \(\cB\sqcup I\), with \(b\sim i\) when
\(i\in I_b\), maps to the incidence graph by \(b\mapsto z_b\).
This map is locally injective: at a block its incident labels are
distinct, and at a label there is at most one incident block over each
point. Thus its girth is also greater than \(g\).

\subsection{The interpolation identity}

For each surviving mark \(a_i=(a_{i,1},\ldots,a_{i,d})\), define the
product Lagrange polynomial
\[
 f_i(X_1,\ldots,X_d)=
 \prod_{\nu=1}^{d}
 \prod_{\substack{\sigma\in S\\\sigma\ne a_{i,\nu}}}
 \frac{X_\nu-\sigma}{a_{i,\nu}-\sigma}\in\F_Q[X_1,\ldots,X_d].
\]
It is \(1\) at \(a_i\) and \(0\) at all other grid points.
Let \(f=(f_i)_{i\in I}\) be the column of these polynomials.
The total degree of every entry of \(ff^{\mathsf T}\) is at most
\[
 2d(|S|-1)<Q-1.
\]

A polynomial in one variable of degree less than \(Q-1\) sums to zero
over \(\F_Q\). For the constant polynomial this is \(Q=0\) in \(\F_Q\).
For \(T^j\), \(1\le j<Q-1\), choose
\(\omega\in\F_Q^\times\) with \(\omega^j\ne1\), which exists because
\(T^j-1\) has at most \(j\) roots; multiplication by \(\omega\)
permutes the inputs and forces the sum to vanish.
Parametrizing an affine line therefore gives
\begin{equation}\label{eq:line-certificate}
 \sum_{z\in\ell_i^\circ}f(z)f(z)^{\mathsf T}=-E_{ii}.
\end{equation}

\begin{lemma}[Rank certificate]\label{lem:certificate}
Let \(D\) be a division ring of characteristic \(p\), and let \(Y_i\)
be arbitrary \(r\times r\) matrices over \(D\). Put
\(Z(z)=\sum_{i\in I(z)}Y_i\). Then
\begin{equation}\label{eq:certificate-rank}
 \sum_{i\in I}\rk_D Y_i\le\sum_{z\in P}\rk_D Z(z).
\end{equation}
\end{lemma}

\begin{proof}
The division ring \(D\) need not contain \(\F_Q\), so we adjoin
\(\F_Q\) inside a matrix ring. Let \(\lambda:\F_Q\to M_a(\F_p)\) be
the regular representation: \(\lambda(c)\) is the matrix of
multiplication by \(c\) in the basis \(1,X,\ldots,X^{a-1}\). It is an
injective ring homomorphism. The prime field \(\F_p\) is central in
\(D\), so in \(M_a(D)\) every \(\lambda(c)\) commutes with every
scalar matrix \(yI_a\), \(y\in D\). Let \(R\subseteq M_a(D)\) be the
subring generated by \(\lambda(\F_Q)\) and \(DI_a\). Then
\(\lambda(\F_Q)\) is a central subfield of \(R\), and
\(y\mapsto yI_a\) embeds \(D\) in \(R\). We regard \(R\) as an
\(\F_Q\)-algebra through \(\lambda\), and matrices over \(D\) as
matrices over \(R\).

Over \(R\) we have
\begin{equation}\label{eq:certificate}
 -\bigoplus_{i\in I}Y_i
 =\sum_{z\in P}(f(z)\otimes I_r)\,Z(z)\,
                 (f(z)^{\mathsf T}\otimes I_r).
\end{equation}
Indeed, expand the right side, exchange the finite sums, and
apply~\eqref{eq:line-certificate} to each label \(i\); the scalars
from \(\F_Q\) are central in \(R\).
This gives \(-\sum_i E_{ii}\otimes Y_i\).

Now read an \(n\times n'\) matrix over \(R\) as an \(na\times n'a\)
matrix over \(D\) made of \(a\times a\) blocks; this respects sums
and products. A matrix \(Y\) over \(D\) becomes \(Y\otimes I_a\),
which a permutation of rows and columns turns into
\(I_a\otimes Y\). Its rank over \(D\) is therefore \(a\rk_D Y\).
Hence the left side of~\eqref{eq:certificate} has rank
\(a\sum_i\rk_D Y_i\), whereas the summand at \(z\) on the right has
rank at most \(a\rk_D Z(z)\). Subadditivity of rank, and division
by \(a\), prove~\eqref{eq:certificate-rank}.
\end{proof}

The identity is purely algebraic. The group in the next section
arranges that each regular \(Z(z)\) is a sum of two matrices of rank
at most one.
